\begin{lemma}
\label{lemma-glue-pseudo-coherent}
Let $R$ be a ring. Let $f_1, \ldots, f_r \in R$ be elements which
generate the unit ideal. Let $m \in \mathbf{Z}$. Let $K^\bullet$
be a complex of $R$-modules. If for each $i$ the complex
$K^\bullet \otimes_R R_{f_i}$ is $m$-pseudo-coherent
(resp.\ pseudo-coherent), then $K^\bullet$ is $m$-pseudo-coherent
(resp.\ pseudo-coherent).
\end{lemma}

\begin{proof}
We will use without further mention that $- \otimes_R R_{f_i}$ is
an exact functor and that therefore
$$
H^i(K^\bullet)_{f_i} =
H^i(K^\bullet) \otimes_R R_{f_i} = H^i(K^\bullet \otimes_R R_{f_i}).
$$
Assume $K^\bullet \otimes_R R_{f_i}$ is $m$-pseudo-coherent
for $i = 1, \ldots, r$. Let $n \in \mathbf{Z}$ be the largest
integer such that $H^n(K^\bullet \otimes_R R_{f_i})$ is nonzero
for some $i$. This implies in particular that $H^i(K^\bullet) = 0$
for $i > n$ (and that $H^n(K^\bullet) \not = 0$) see
Algebra, Lemma \ref{algebra-lemma-cover}.
We will prove the lemma by induction on $n - m$.
If $n < m$, then the lemma is true by
Lemma \ref{lemma-recognize-pseudo-coherent}.
If $n \geq m$, then $H^n(K^\bullet)_{f_i}$ is a finite $R_{f_i}$-module
for each $i$, see
Lemma \ref{lemma-finite-cohomology}.
Hence $H^n(K^\bullet)$ is a finite $R$-module, see
Algebra, Lemma \ref{algebra-lemma-cover}.
Choose a finite free $R$-module $E$ and a surjection $E \to H^n(K^\bullet)$.
As $E$ is projective we can lift this to a map of complexes
$\alpha : E[-n] \to K^\bullet$. Then the cone $C(\alpha)^\bullet$ has
vanishing cohomology in degrees $\geq n$. On the other hand, the
complexes $C(\alpha)^\bullet \otimes_R R_{f_i}$ are $m$-pseudo-coherent
for each $i$, see
Lemma \ref{lemma-cone-pseudo-coherent}.
Hence by induction we see that $C(\alpha)^\bullet$ is $m$-pseudo-coherent
as a complex of $R$-modules. Applying
Lemma \ref{lemma-cone-pseudo-coherent}
once more we conclude.
\end{proof}
